\setlength{\chnumsep}{12em}
\chapter{Sequences}

\section{Preliminaries}

Our aim is here to rigorously deal with the various results you have probably already encountered
in calculus. While we will be \emph{rigorous}, we won't be \emph{abstract}. You have
probably already encountered limits of functions. One way to think about
such limits is \emph{sequences}. In highschool mathematics,
you might have encountered ``intuitive'' notions of sequences:
\[
\text{Sequence: }  \text{a sequence is a list of numbers, e.g. } a_1, a_2, a_3, \ldots
\]
Let us work with this for now. Consider the following limit:
\[
\lim_{x \to 0} \sin(x) = 0.
\]
One way to think about it is by getting closer and closer to $0$ in a continuous manner. However, an equivalent,
and maybe more powerful way is to think about a list of numbers that grows closer
and closer to $0$. So, let us plot $\sin(x)$ for $x = x_1, x_2, x_3, \ldots$ where
$x_n \to 0$. This allows us to see that $\sin(x)$ approaches $0$ as $x$ grows closer to $0$.

\begin{figure}[H]
\begin{fullwidth}
\centering
\begin{tikzpicture}
\begin{axis}[
   axis lines = middle,
   xlabel = {$x$},
   ylabel = {$\sin(x)$},
   xlabel style={anchor=west},
   ylabel style={anchor=south},
   xmin=-0.3, xmax=2.4,
   ymin=-0.6, ymax=1.25,
   xtick={0,0.5,1,1.5,2},
   ytick={-1,-0.5,0,0.5,1},
   samples=200,
   domain=-0.3:2.4,
   width=15cm, height=7cm,
   scale only axis,
   clip=false,
]

% the curve y = sin(x), x in radians
\addplot[thick, blue] {sin(deg(x))};

% the sequence x_i = 2/(i+1) -> 0, and its images under sin
   \draw[red!70, dashed] (axis cs:2.0000,0) -- (axis cs:2.0000,0.9093);
   \filldraw[red] (axis cs:2.0000,0) circle (1.5pt);
   \filldraw[red] (axis cs:2.0000,0.9093) circle (1.5pt);
   \draw[red!70, dashed] (axis cs:1.0000,0) -- (axis cs:1.0000,0.8415);
   \filldraw[red] (axis cs:1.0000,0) circle (1.5pt);
   \filldraw[red] (axis cs:1.0000,0.8415) circle (1.5pt);
   \draw[red!70, dashed] (axis cs:0.6667,0) -- (axis cs:0.6667,0.6184);
   \filldraw[red] (axis cs:0.6667,0) circle (1.5pt);
   \filldraw[red] (axis cs:0.6667,0.6184) circle (1.5pt);
   \draw[red!70, dashed] (axis cs:0.5000,0) -- (axis cs:0.5000,0.4794);
   \filldraw[red] (axis cs:0.5000,0) circle (1.5pt);
   \filldraw[red] (axis cs:0.5000,0.4794) circle (1.5pt);
   \draw[red!70, dashed] (axis cs:0.4000,0) -- (axis cs:0.4000,0.3894);
   \filldraw[red] (axis cs:0.4000,0) circle (1.5pt);
   \filldraw[red] (axis cs:0.4000,0.3894) circle (1.5pt);
   \draw[red!70, dashed] (axis cs:0.3333,0) -- (axis cs:0.3333,0.3272);
   \filldraw[red] (axis cs:0.3333,0) circle (1.5pt);
   \filldraw[red] (axis cs:0.3333,0.3272) circle (1.5pt);
   \draw[red!70, dashed] (axis cs:0.2857,0) -- (axis cs:0.2857,0.2818);
   \filldraw[red] (axis cs:0.2857,0) circle (1.5pt);
   \filldraw[red] (axis cs:0.2857,0.2818) circle (1.5pt);
   \draw[red!70, dashed] (axis cs:0.2000,0) -- (axis cs:0.2000,0.1987);
   \filldraw[red] (axis cs:0.2000,0) circle (1.5pt);
   \filldraw[red] (axis cs:0.2000,0.1987) circle (1.5pt);
   \draw[red!70, dashed] (axis cs:0.1538,0) -- (axis cs:0.1538,0.1533);
   \filldraw[red] (axis cs:0.1538,0) circle (1.5pt);
   \filldraw[red] (axis cs:0.1538,0.1533) circle (1.5pt);
   \draw[red!70, dashed] (axis cs:0.1111,0) -- (axis cs:0.1111,0.1109);
   \filldraw[red] (axis cs:0.1111,0) circle (1.5pt);
   \filldraw[red] (axis cs:0.1111,0.1109) circle (1.5pt);
   \draw[red!70, dashed] (axis cs:0.0800,0) -- (axis cs:0.0800,0.0799);
   \filldraw[red] (axis cs:0.0800,0) circle (1.5pt);
   \filldraw[red] (axis cs:0.0800,0.0799) circle (1.5pt);
   \draw[red!70, dashed] (axis cs:0.0606,0) -- (axis cs:0.0606,0.0606);
   \filldraw[red] (axis cs:0.0606,0) circle (1.5pt);
   \filldraw[red] (axis cs:0.0606,0.0606) circle (1.5pt);
   \draw[red!70, dashed] (axis cs:0.0455,0) -- (axis cs:0.0455,0.0454);
   \filldraw[red] (axis cs:0.0455,0) circle (1.5pt);
   \filldraw[red] (axis cs:0.0455,0.0454) circle (1.5pt);
   \draw[red!70, dashed] (axis cs:0.0333,0) -- (axis cs:0.0333,0.0333);
   \filldraw[red] (axis cs:0.0333,0) circle (1.5pt);
   \filldraw[red] (axis cs:0.0333,0.0333) circle (1.5pt);
   \draw[red!70, dashed] (axis cs:0.0241,0) -- (axis cs:0.0241,0.0241);
   \filldraw[red] (axis cs:0.0241,0) circle (1.5pt);
   \filldraw[red] (axis cs:0.0241,0.0241) circle (1.5pt);
   \draw[red!70, dashed] (axis cs:0.0171,0) -- (axis cs:0.0171,0.0171);
   \filldraw[red] (axis cs:0.0171,0) circle (1.5pt);
   \filldraw[red] (axis cs:0.0171,0.0171) circle (1.5pt);
   \draw[red!70, dashed] (axis cs:0.0121,0) -- (axis cs:0.0121,0.0121);
   \filldraw[red] (axis cs:0.0121,0) circle (1.5pt);
   \filldraw[red] (axis cs:0.0121,0.0121) circle (1.5pt);
% first few terms labelled
\node[red, font=\footnotesize] at (axis cs:2.0000,-0.18) {$x_0$};
\node[red, font=\footnotesize] at (axis cs:1.0000,-0.18) {$x_1$};
\node[red, font=\footnotesize] at (axis cs:0.6667,-0.18) {$x_2$};
\node[red, font=\footnotesize] at (axis cs:0.5000,-0.18) {$x_3$};

% origin
\filldraw[black] (axis cs:0,0) circle (1.5pt);
\node[below left, font=\footnotesize] at (axis cs:0,0) {$0$};

% sequence formula

% continuation / convergence note
\node[red!70, font=\footnotesize] at (axis cs:0.55,-0.34) {$x_4,\, x_5,\, \ldots \to 0$};
\draw[->, thick, red!70] (axis cs:0.55,-0.28) -- (axis cs:0.18,-0.1);

\end{axis}
\end{tikzpicture}
\label{fig:sin-sequence}
\end{fullwidth}
\vspace{-\baselineskip}\vspace{-\baselineskip}
  \sideparmargin{inner}
  \sidepar{\vspace{\baselineskip}
    \caption{A sequence $x_i \to 0$ and the corresponding values $\sin(x_i) \to 0$.}}
\end{figure}

This is primarily the reason we work with sequences. Although it is easy to imagine the limit
in the other way too, it is often helpful to work with a discrete set of points instead.

However the definition of sequence we laid out hardly suffices as a good definition, for one, we have postponed the burden of definition
from sequence to a list of numbers. Let us give a more formal definition.

\begin{definition}
    A \emph{sequence} is a function $a: \mathbb{N} \to \RR$, such that
    \[
        a_n = a(n)
    \]
    for all $n \in \mathbb{N}$ is the $n$-th term of the sequence.
\end{definition}

\marginnote{I will adopt the convention that $\mathbb{N}$ does not include $0$. The notions of functions
and other primitive things will be assumed. Sequences will be denoted by $(a_n)_{n=1}^{\infty}$ or
simply $(a_n)$, for reasons that will become clear as we go.}

\begin{example}
    Some examples from highschool mathematics include the sequences defined by $a_n = ar^n$,
    and $a_n = a + nd$. These sequences are \emph{geometric} and \emph{arithmetic}, respectively.
\end{example}


For this course we will assume the properties and existence of $\RR$. In more rigorous analysis
courses, you will learn about the \emph{construction} of $\RR$, but we will not go into detail
here. The properties of $\RR$ are listed below.

\begin{itemize}
    \item For all $a, b \in \RR$, $a + b, a - b, a \cdot b, a / b$ are defined (for the last, $b \neq 0$).
    \item For all $a, b \in \RR$, either $a < b$, $a = b$, or $a > b$.
\end{itemize}

You might realise that these properties are also the properties that $\QQ$ has.
What, then, distinguishes $\QQ$ from $\RR$? Clearly they're not the same, we've all seen
proofs that $\sqrt{2}$ is in fact not rational, but we also happen to ``know'' that $\sqrt{2}$ is a
real number. There is precisely one property that distinguishes $\QQ$ from $\RR$, and it happens
to unfortunately be a little more obtuse than the above, rather obvious ones.

The property we speak of is the \emph{Least Upper Bound Property}, but before stating it,
we need to define what an \emph{upper bound} is.

\begin{definition}
    A \emph{least upper bound} of a set $S$ is a real number
    $b$ such that $b \geq a$ for all $a \in S$.
\end{definition}

We now must define what it means to have a \emph{least upper bound} for a set $S$:
\begin{definition}
    A set $S$ has a \emph{least upper bound} if there exists a real number $b$ such that
    $b$ is an upper bound for $S$, and if $a$ is an upper bound for $S$, then $b \geq a$. We
    write $b = \sup S$.
\end{definition}

Now we can actually state this mysterious property of the reals:
\[
    \text{If } S \subset \RR \text{ has an upper bound } b, \text{ then } S \text{ has a least upper bound } b.
\]

There is a similar property for lower bounds, that I will leave for you to prove, using
the least upper bound property.
\begin{definition}
    A set $S$ has a \emph{least lower bound} if there exists a real number $b$ such that
    $b$ is a lower bound for $S$, and if $a$ is a lower bound for $S$, then $b \leq a$. We
    write $b = \inf S$.
\end{definition}

One final property that we require is the archimedean property. This seems like a
really obvious result, so much so that you might be surprised that it is even necessary.
However as we will continue to do in the course, we will not take it as a given
and actually prove it.

\begin{theorem}[Archimedean Property]
	For any real number $a$, there exists a positive integer $N$ such that $N > a$.
\end{theorem}

\begin{proof}
	Suppose no such $N$ exists. Then for every positive integer $N$, $N \leq a$,
	so $a$ is an upper bound of $\NN \subset \RR$. By the Least Upper Bound Property,
	there must exist a least upper bound $b$ of $\NN \subset \RR$.

	Note that $b-1$ must be an upper bound of $\NN$ as well, because if $b - 1 < n \in \NN$,
	then $b < n + 1 \in \NN$, which cannot be true since $b$ is an upper bound of $\NN$.
	However, if $b-1$ is an upper bound of $\NN$, then $b - 1 > b$ by the definition
	of Least Upper Bound, which cannot be true! Hence some $N \in \NN$ must exist such that $N > a$.
\end{proof}

\begin{exc}
	\begin{exercise}
		Prove that if a set has a lower bound, then it has a greatest lower bound.
	\end{exercise}
\begin{solution}

		Let $S \subseteq \mathbb{R}$ be nonempty and bounded below. Define
		\[
			T = \{-s : s \in S\}.
		\]
		Since $S$ is bounded below, there exists $m \in \mathbb{R}$ such that $m \le s$ for all $s \in S$. Then $-s \le -m$ for all $s \in S$, so $T$ is bounded above by $-m$.

		By the least upper bound property of $\mathbb{R}$, $T$ has a supremum; call it $\alpha = \sup T$.

		\begin{claim}
			$-\alpha = \inf S$.
		\end{claim}

		\emph{Lower bound.} For all $t \in T$, $t \le \alpha$, so for all $s \in S$, $-s \le \alpha$, i.e.\ $-\alpha \le s$. Thus $-\alpha$ is a lower bound for $S$.

		\emph{Greatest lower bound.} Let $\beta$ be any lower bound for $S$, so $\beta \le s$ for all $s \in S$. Then $-s \le -\beta$ for all $s \in S$, i.e.\ $-\beta$ is an upper bound for $T$. Since $\alpha = \sup T$ is the \emph{least} upper bound, $\alpha \le -\beta$, so $\beta \le -\alpha$.

		Hence $-\alpha$ is a lower bound for $S$ and is $\ge$ every other lower bound, so $-\alpha = \inf S$.
\end{solution}
\end{exc}

\section{Convergent Sequences}

Although just about anything can be defined as a sequence---for example the oscillating sequence
defined by $a_n = (-1)^n$---our interest resides in a more limited set of ``nice'' sequences. Consider
for example the sequence defined by $a_n = 1 + 1/n$. As we go to farther values of $n$, this
sequence grows closer and closer to the value $1$. This notion of growing closer and closer to a value
is captured by the concept of convergence (which is equivalent to the idea of a limit).

\begin{definition}
    A sequence $(a_n)$ is \emph{convergent}, and is said to converge to $\limsup a_i$, called the
   \vocab{limit} of the sequence, if for every $\varepsilon > 0$, there exists an $N \in \NN$ such that
    $\abs{a_n - \limsup a_i} < \varepsilon$ for all $n \geq N$, where $\limsup a_i$ is the limit of the sequence.
\end{definition}

We write $\lim_{n \to \infty} a_n = \limsup a_i$ to denote that the sequence $(a_n)$ converges to $\limsup a_i$.
If the definition looks a little obtuse, then do not worry. At first it might seem terribly confusing,
but this actually just captures the notion we just mentioned, the point is that by going to large enough
values, we can get arbitrarily close to the value it converges to.

% Preamble requirements:
% \usepackage{tikz}
% \usetikzlibrary{decorations.pathreplacing}
% (assumes tufte-latex or similar package providing \marginfigure)

\begin{marginfigure}
\centering
\begin{tikzpicture}[scale=0.82, every node/.style={font=\scriptsize}]

  \def\dr{1pt} % dot radius

  % main axis
  \draw[->] (-3.5,0) -- (2.4,0);

  % isolated points a_1, a_2, a_3
  \foreach \x/\lbl in {-3.3/a_1, -2.85/a_2, -2.4/a_3} {
    \fill (\x,0) circle (\dr);
    \node[above] at (\x,0.08) {$\lbl$};
  }

  % trailing dots, sitting on the axis like the other points
  \foreach \x in {-2.05,-1.92,-1.79} {
    \fill (\x,0) circle (1pt);
  }

  % interval endpoints
  \def\L{-1.5}
  \def\R{1.5}
  \def\A{0.5}

  % open interval parentheses
  \node[scale=1.3] at (\L,0) {$($};
  \node[scale=1.3] at (\R,0) {$)$};

  % a_N: first point inside the interval, near the left end
  \fill (-1.3,0) circle (\dr);

  % dense, congested cluster of points near a
  \foreach \x in {0.0,0.08,0.17,0.4,0.55,0.8,0.95} {
    \fill (\x,0) circle (\dr);
  }

  % labels below the axis
  \node[below] at (\L,-0.16) {$a-\varepsilon$};
  \node[below] at (\A,-0.16) {$a$};
  \node[below] at (\R,-0.16) {$a+\varepsilon$};

  % a_N label with downward arrow
  \node at (-1.3,0.7) {$a_N$};
  \draw[->] (-1.3,0.55) -- (-1.3,0.1);

  % brace for V_epsilon(a)
  \draw[decorate,decoration={brace,amplitude=4pt}]
    (\L,1.0) -- (\R,1.0) node[midway,above=5pt] {$V_\varepsilon(a)$};

\end{tikzpicture}
\caption{A sequence that converges to $a$}
\end{marginfigure}

\begin{remark}
    There is a \emph{topological} way to think about this. Unfortunately the course won't go
    into this in detail, but something useful is the notion of $\varepsilon$ neighbourhood of a point.
    This is simply the set of all points within $\varepsilon$ distance of the given point. So how
    is it different from the definition we just gave?

    The notion of distance is not just limited to $\RR$---while we do have $\abs{x-a}$ as
    the distance of $x$ from $a$, we can also find out the distance between two points in for
    example $\RR^2$, using what is called the \emph{metric}. For now these are just words without
    meaning, but perhaps the idea of convergence by being able to find a point in the sequence for
    any $\epsilon > 0$, such that all terms after it are within the $\varepsilon$ neighbourhood,
    $V_{\varepsilon}$ of that point.
\end{remark}

A good imperative to have when defining something is one, to actually see if this object exists,
and two to see that it is well defined. When we write down a symbol like $\lim_{n \to \infty} a_n$,
we must check if this is actually a unique value, that is there is only value that the sequence
converges to. Otherwise our definition would be ambiguous, the symbol would mean too many things.
Therefore, let us do a little proposition.

\begin{proposition*}
    If the sequence $(a_n)$ converges to $\limsup a_i$, and $M$, then $\limsup a_i = M$.
\end{proposition*}

\begin{proof}
	Suppose the sequence $(a_n)$ converges to $\limsup a_i$, and $M$. Then, by the
	definition of convergence, for any $\varepsilon > 0$, there exists $N_1$ such that for all $n \geq N_1$, $\abs{a_n - \limsup a_i} < \varepsilon$.
	Similarly, there exists $N_2$ such that for all $n \geq N_2$, $\abs{a_n - M} < \varepsilon$.

	A straightforward way to proceed is to add up the two inequalities. But first
	you must remember that the $n$ that satisfied $n > N_1$ does not necessarily satisfy $n > N_2$.
	How do we deal with this? Pick any $n > \max(N_1, N_2)$, so that now $a_n$ satisfies
	both the inequalities.
	\[
	\abs{a_n - \limsup a_i} + \abs{a_n - M} = \abs{a_n - \limsup a_i} + \abs{M - a_n} < 2\varepsilon,
	\]
	and now we use one of the most important inequalities in these epsilon proofs, the triangle inequality:
	\[
	2\varepsilon > \abs{a_n - \limsup a_i} + \abs{a_n - M} \ge \abs{a_n - \limsup a_i} + \abs{M - a_n} = \abs{M - \limsup a_i}.
	\]

	But remember that we could choose \emph{any} epsilon here. Now there are two ways to proceed.
	One, by choosing the initial
	epsilon to be $\varepsilon'/2$, we have that for any $\varepsilon' > 0$,
	\[
	\varepsilon' > \abs{M - \limsup a_i}
	\]
	but that must mean $M - \limsup a_i = 0$! (convince yourself of it, by arguing formally using
	$M-\limsup a_i = r \neq 0$).

	A more clever proof would be to use $\varepsilon = \abs{M - \limsup a_i}/2$, assuming for contradiction
	that $M \neq \limsup a_i$, so that we would get:
	\[
	 \abs{M - \limsup a_i} < 2\varepsilon = \abs{M - \limsup a_i}
	\]
	which cannot be! The idea behind this is to think about close $a_n$ must be to
	$\limsup a_i$ and $M$, and achieve a contradiction by thinking about the distance between
	the two values.
\end{proof}

Now we know that our symbols are well defined. But I have yet to show you an example
of a sequence that converges to some value using the definition we adopted. Let us
now do that.

\begin{example}
The sequence $a_n = \dfrac{1}{n}$ converges to $0$.

\hrulebar

Before writing the proof, let us think about what
we need. The definition requires: for any $\varepsilon > 0$, there exists
$N \in \mathbb{N}$ such that $\abs{a_n - 0} < \varepsilon$ for all $n > N$. The
quantifier, for any, matters here. $\varepsilon$ is given to us first (adversarially, we may
imagine), and only afterward do we get to choose $N$; the choice of $N$ is
therefore allowed to depend on $\varepsilon$.

Fix an $\varepsilon > 0$, how large must $n$ be to guarantee
$\abs{1/n - 0} < \varepsilon$? Since $1/n > 0$, this inequality is equivalent to
$n > 1/\varepsilon$. This tells us exactly how to choose $N$: any natural
number $N \geq 1/\varepsilon$ will do. We now write this up formally.

\begin{proof} Let $\varepsilon > 0$. By the Archimedean property,
there exists $N \in \mathbb{N}$ with $N > 1/\varepsilon$. We claim this $N$
works. Let $n > N$. Then
\[
n > N > \frac{1}{\varepsilon} \quad\Longrightarrow\quad \frac{1}{n} < \varepsilon,
\]
and since $n \geq 1$ we have $1/n > 0$, so
\[
\abs{a_n - 0} = \frac{1}{n} < \varepsilon.
\]
As $\varepsilon > 0$ was arbitrary, this shows $a_n \to 0$. $\blacksquare$
\end{proof}

\end{example}


\begin{example}
The sequence $a_n = \dfrac{2n}{3n+1}$ converges to $\dfrac{2}{3}$.

\hrulebar

As before, we first manipulate $\abs{a_n - \limsup a_i}$
algebraically until the dependence on $n$ becomes transparent:
\[
\abs{\frac{2n}{3n+1} - \frac{2}{3}}
= \abs{\frac{3(2n) - 2(3n+1)}{3(3n+1)}}
= \abs{\frac{-2}{3(3n+1)}}
= \frac{2}{3(3n+1)}.
\]
This is a decreasing function of $n$, so making it small is just a matter of
making $n$ large. We want
\[
\frac{2}{3(3n+1)} < \varepsilon
\quad\Longleftrightarrow\quad
3n + 1 > \frac{2}{3\varepsilon}
\quad\Longleftarrow\quad
n > \frac{2}{9\varepsilon}.
\]
(We used $\Longleftarrow$ rather than $\Longleftrightarrow$ in the last step
purely for convenience — we don't need the sharpest possible $N$, only
\emph{some} $N$ that works, so we are free to bound things generously.) This
suggests taking $N$ to be any natural number exceeding $2/(9\varepsilon)$.

\medskip
\noindent\textbf{Proof.} Let $\varepsilon > 0$, and choose $N \in \mathbb{N}$
with $N > \dfrac{2}{9\varepsilon}$ (possible by the Archimedean property). Let
$n > N$. Then $n > 2/(9\varepsilon)$, so $9\varepsilon n > 2$, and hence
$3(3n+1) = 9n + 3 > 9n > \dfrac{2}{\varepsilon}$. Rearranging,
\[
\frac{2}{3(3n+1)} < \varepsilon.
\]
Combined with the algebraic identity from the discussion above,
\[
\abs{\frac{2n}{3n+1} - \frac{2}{3}} = \frac{2}{3(3n+1)} < \varepsilon.
\]
Since $\varepsilon > 0$ was arbitrary, $a_n \to 2/3$. $\blacksquare$
\end{example}

You might've realised that this is often quite cumbersome to work with. For example,
what about proving $a_n = 1/n + 2n/(3n + 1)$ converged to $2/3$? Try doing this yourself,
but be warned, this will be annoying to deal with.

We can simplify many matters here by proving some useful properties of limits of
sequences. These are set up as exercises to be done at the end of this section. They're
guided of course, since picking the right epsilon is a skill that requires learning
which I will not presume.

\begin{remark}
    A subsequent definition that we can make is to formalise the notion of divergence. A
    sequence $(a_n)$ is said to \emph{diverge} if there does not exist a value that the sequence
    converges to.

    However it serves to distinguist between two different notions of divergence. Consider the
    sequences $a_n = (-1)^n$ and $b_n = n$. Both sequences diverge, but the first one
    has no well defined behaviour as we grow larger, while the larger gets arbitrarily large.

    To formalise this distinction, we often work in the \emph{extended real numbers} defined by
    $\RR \union \{\infty, -\infty\}$, where $\infty$ and $-\infty$
    represent the notions of positive and negative infinity, respectively. We define a few properties:

    \begin{itemize}
        \item For any $a \in \RR$, $a \leq \infty$ and $a \geq -\infty$.
        \item For any $a \in \RR$, $a + \infty = \infty$ and $- \infty + a = -\infty$.
    \end{itemize}

    Now we are able to extend the notion of convergence:

    \begin{itemize}
        \item A sequence $(a_n)$ converges to $\infty$ if for any $M \in \RR$, there exists a $N$ such that $a_n > M$ for all $n \geq N$.
        \item A sequence $(a_n)$ converges to $-\infty$ if for any $M \in \RR$, there exists a $N$ such that $a_n < M$ for all $n \geq N$.

    \end{itemize}
    \label{rem:extended_convergence}
\end{remark}

\begin{exc}

    We will prove the following properties of limits of sequences:
    \begin{itemize}
        \item If $\lim_{n \to \infty} a_n = a$, and $\lim_{n \to \infty} b_n = b$, and $c_n$ is the sequence defined by $c_n = a_n + b_n$, then $\lim_{n \to \infty} c_n = a + b$.
        \item If $\lim_{n \to \infty} a_n = a$, and $\lim_{n \to \infty} b_n = b$, and $c_n$ is the sequence defined by $c_n = a_n \cdot b_n$, then $\lim_{n \to \infty} c_n = a \cdot b$.
        \item If $\lim_{n \to \infty} a_n = a$, and $\lim_{n \to \infty} b_n = b \neq 0$, and $c_n$ is the sequence defined by $c_n = {a_n}/{b_n}$ when $b_n \neq 0$, and $0$ when $b_n = 0$, then $\lim_{n \to \infty} c_n = {a}/{b}$.
    \end{itemize}
    \begin{exercise}
    Prove that if
    \[
    \abs{x - x_0} < \frac{\varepsilon}{2} \quad \text{and} \quad \abs{y - y_0} < \frac{\varepsilon}{2},
    \]
    then
    \[
    \abs{(x+y) - (x_0+y_0)} < \varepsilon,
    \]
    \[
    \abs{(x-y) - (x_0-y_0)} < \varepsilon.
    \]
    \end{exercise}

    \begin{exercise}
    Prove that if
    \[
    \abs{x - x_0} < \min\left( \frac{\varepsilon}{2(\abs{y_0}+1)}, 1 \right)
    \quad \text{and} \quad
    \abs{y - y_0} < \frac{\varepsilon}{2(\abs{x_0}+1)},
    \]
    then $\abs{xy - x_0y_0} < \varepsilon$.

    (The notation ``min'' is self evident, the first inequality in the
    hypothesis just means that
    \[
    \abs{x - x_0} < \frac{\varepsilon}{2(\abs{y_0}+1)} \quad \text{and} \quad \abs{x - x_0} < 1;
    \]
    at one point in the argument you will need the first inequality, and at
    another point you will need the second. One more word of advice: since the
    hypotheses only provide information about $x - x_0$ and $y - y_0$, it is
    almost a foregone conclusion that the proof will depend upon writing $xy -
    x_0y_0$ in a way that involves $x - x_0$ and $y - y_0$.)
    \end{exercise}

    \begin{exercise}
    Prove that if $y_0 \neq 0$ and
    \[
    \abs{y - y_0} < \min\left( \frac{\abs{y_0}}{2}, \frac{\varepsilon\abs{y_0}^2}{2} \right),
    \]
    then $y \neq 0$ and
    \[
    \abs{\frac{1}{y} - \frac{1}{y_0}} < \varepsilon.
    \]
    \end{exercise}

    \begin{exercise}
	Now prove the three propositions we stated at the beginning.
	\label{exc:limit_properties}
    \end{exercise}
   	\begin{solution}
		\textbf{(i) Sum of limits}
		\begin{claim}
			\[\lim_{n\to\infty}(a_n+b_n) = \lim_{n\to\infty} a_n + \lim_{n\to\infty} b_n.\]
		\end{claim}
		\begin{proof}
			Let $\varepsilon > 0$. Since $a_n \to a$ and $b_n \to b$, choose $N_a, N_b \in \mathbb{N}$ such that
			\[
				n \geq N_a \implies \abs{a_n - a} < \frac{\varepsilon}{2}, \qquad
				n \geq N_b \implies \abs{b_n - b} < \frac{\varepsilon}{2}.
			\]
			Let $N = \max(N_a, N_b)$. Then for all $n \geq N$,
			\[
				\abs{(a_n+b_n) - (a+b)} \leq \abs{a_n - a} + \abs{b_n - b} < \varepsilon.
			\]
		\end{proof}
		\textbf{(ii) Product of limits}
		\begin{claim}
			\[\lim_{n\to\infty} (a_n b_n) = (\lim_{n\to\infty} a_n)(\lim_{n\to\infty} b_n).\]
		\end{claim}
		\begin{proof}
			Let $\varepsilon > 0$. Choose $N_a, N_b \in \mathbb{N}$ such that
			\begin{align*}
				n \geq N_a &\implies \abs{a_n - a} < \frac{\varepsilon}{2(3\abs{b}+1)}, \\
				n \geq N_b &\implies \abs{b_n - b} < \min\left(\frac{\varepsilon}{2(\abs{a}+1)}, \frac{\abs{b}}{2}\right).
			\end{align*}
			For $n \geq N_b$, the reverse triangle inequality gives $\abs{b} - \abs{b_n} \leq \abs{b - b_n} < \abs{b}/2$, so $\abs{b_n} < \tfrac{3}{2}\abs{b}$.

			Then for $n \geq N := \max(N_a, N_b)$,
			\begin{align*}
				\abs{a_nb_n - ab} &= \abs{b_n(a_n-a) + a(b_n-b)} \\
				&\leq \abs{b_n}\,\abs{a_n - a} + \abs{a}\abs{b_n - b} \\
				&< \frac{3\abs{b}}{2}\cdot\frac{\varepsilon}{2(3\abs{b}+1)} + \frac{3\abs{a}}{2}\cdot\frac{\varepsilon}{2(3\abs{a}+1)} \\
				&< \frac{\varepsilon}{2} + \frac{\varepsilon}{2} = \varepsilon.
			\end{align*}
		\end{proof}
		\textbf{(iii) Quotient of limits}
		\begin{claim}
			If $\lim_{n\to\infty} a_n = a$ and $\lim_{n\to\infty} b_n = b$, then
			\[
				\lim_{n\to\infty} \frac{a_n}{b_n} \mathbb{1}_{b_n\neq 0} = \frac{a}{b}.
			\]
		\end{claim}
		\begin{proof}
			First, we show that
			\[
				\frac{1}{b_n} \to \frac{1}{b}.
			\]
			Since $b_n \to b$ and $b \neq 0$, there exists $N_1 \in \mathbb{N}$ such that
			\[
				n \geq N_1 \implies \abs{b_n - b} < \frac{\abs{b}}{2}.
			\]
			Hence, for $n \geq N_1$,
			\[
				\abs{b_n} \geq \abs{b} - \abs{b_n - b} > \frac{\abs{b}}{2},
			\]
			so in particular $b_n \neq 0$.

			Now let $\varepsilon>0$. Since $b_n\to b$, choose $N_2\in\mathbb{N}$ such that
			\[
				n\geq N_2
				\implies
				\abs{b_n - b} < \frac{\abs{b}}{2}.
			\]
			Hence, for $n \geq N_2$,
			\[
				\abs{b_n} \geq \abs{b} - \abs{b_n - b} > \frac{\abs{b}}{2},
			\]
			so in particular $b_n \neq 0$.

			Now let $\varepsilon>0$. Since $b_n\to b$, choose $N_2\in\mathbb{N}$ such that
			\[
				n\geq N_2
				\implies
				\abs{b_n - b} < \frac{\abs{b}}{2}.
			\]
			Thus, for $n\geq N:=\max(N_1,N_2)$,
			\[
				\abs{\frac{1}{b_n}-\frac{1}{b}}
				=
				\frac{\abs{b_n-b}}{\abs{b_n}\,\abs{b}}
				<
				\frac{\varepsilon \abs{b}^2/2}{(b/2)\abs{b}}
				=
				\varepsilon.
			\]
			Therefore,
			\[
				\frac{1}{b_n}\to\frac{1}{b}.
			\]

			By part (ii), the product of convergent sequences converges to
			the product of their limits. Hence
			\[
				\frac{a_n}{b_n}
				=
				a_n\frac{1}{b_n}
				\to
				a\frac{1}{b}
				=
				\frac{a}{b}.
			\]
		\end{proof}
	\end{solution}

    \begin{exercise}
	We call a sequence bounded iff there exists a number $M$ such that $\abs{x} \leq M$ for all $x$ in the sequence.
	We call a sequence monotonically increasing iff for all $n \leq m$, $x_n \leq x_m$.

	Suppose that $(a_n)$ is a monotonically increasing bounded sequence. Prove that
	it converges.

	\emph{Hint: Let $B = \{a_n : n \in \mathbb{N}\}$, the set of all values of $a_n$. Note that this
	is bounded. Then consider the least upper bound of $B$. How does this relate to the convergence of $(a_n)$?}
	\label{exc:monotone_convergence}
	\end{exercise}
	\begin{solution}
	    Suppose that $(a_n)_{n=1}^{\infty}$ is monotonically increasing. Let
	    $a = \sup\{a_n \mid n \in \mathbb{N}\}$.
	    Then,
	    \[ \lim_{n\to\infty} a_n = a.\]

	    \begin{proof}
	    We proceed by contradiction. Suppose not. Then there exists $\varepsilon > 0$
	    such that
	    \[
	    \nexists N \in \mathbb{N} \text{ so that } \forall n > N,\ \abs{a - a_n} < \varepsilon.
	    \]

	    Note that $a \geq a_n$ for any $n \in \mathbb{N}$ by definition of supremum.

	    \begin{claim}
	    For all $K \in \mathbb{N}$ we have that $a - a_K \geq \varepsilon.$
	    \end{claim}

	    This is true because we know that there must exist $k^* \geq K$ such that $a - a_{k^*} \geq \varepsilon$.
	    If no such $k^*$ existed, our assumption above would not hold and we would be done.

	    Now, since $(a_n)_{n=1}^{\infty}$ is monotonically increasing,
	    \[
	    a_{k^*} \geq a_K \implies a - a_K \geq a - a_{k^*} \geq \varepsilon.
	    \]
	    Thus, $a - a_K \geq \varepsilon$.

	    Now, since $\forall K \in \mathbb{N},\ a - a_K \geq \varepsilon$, we have
	    \[
	    \forall K \in \mathbb{N},\ a - \varepsilon \geq a_K.
	    \]
	    But then $a - \varepsilon$ is an upper bound of $(a_n)_{n=1}^{\infty}$, which cannot
	    be the case, since $a - \varepsilon < a$, violating the definition of $a$ being a
	    supremum. This is a contradiction, so the claim holds.
	    \end{proof}
	\end{solution}

	\begin{exercise}
	Show that every convergent sequence is bounded.
	\end{exercise}
	\begin{solution}
		The sketch of this is as follows. If $(a_n)_{n=1}^{\infty}$ converges to
		$L \in \RR$, then there exist a $N \in \NN$ such that $\abs{a_n - L} \leq 1$ for all
		$n \geq N$. Thus, $\abs{a_n} \leq \abs{L} + 1$ for all $n \geq N$, so $(a_n)_{n=1}^{\infty}$ is bounded
		by $\max(a_1, a_2, \dots, a_N, \abs{L} + 1)$
	\end{solution}
\end{exc}


\section{Cauchy Sequences}

Sometimes we know that a sequence does converge, but using the above definition of convergence
makes it hard to prove that it does. A more powerful definition is the Cauchy sequence definition.

Let us go through an example.
To estimate square roots, a method that has been used since
ancient times is the Babylonian method, which starts with a guess and
repeatedly averages it with $37$ divided by itself. Let us define a sequence using it,
with initial value $6$ (since it is close to the square root of $37$):
\[
x_1 = 6, \qquad x_{n+1} = \frac{1}{2}\left(x_n + \frac{37}{x_n}\right)
\]


\begin{table}[H]
\begin{sidecaption}{The first few terms of the Babylonian square root sequence. $x_4$ and
$x_5$ are so close that they have the same values upto 6 decimal places.}
\begin{center}
\begin{tabular}{c c c}
\toprule
$n$ & $x_n$ & $\abs{x_n - x_{n-1}}$ \\
\midrule
1 & 6.000000 & --- \\
2 & 6.083333 & 0.083333 \\
3 & 6.082763 & 0.000571 \\
4 & 6.082763 & 0.000000 \\
5 & 6.082763 & 0.000000 \\
\bottomrule
\end{tabular}
\end{center}
\end{sidecaption}
\end{table}

I hope it is clear to you that this sequence converges to whatever the value of $\sqrt{37}$ is.
If you're not convinced, write down a few more values of the sequence, and use whatever value
of $\sqrt{37}$ your calculator gives you. But how do you go about proving this?
It seems very inelegant to have to rely on a calculator to prove that a sequence converges,
nevermind if it even suffices as a proof \footnote[-2.2cm]{It can be made rigorous by proving the
decimal expansion of $\sqrt{37}$ to some number of decimal places. Then you could try to show
that a term of the sequence has the same value as the decimal expansion of $\sqrt{37}$ up to that point,
so that their difference is less than $10^{-m}$ for some $m$. I hope that this comes
across as really cumbersome to you.}.

But I have deliberately written down the third column of the table here,
so that you can see that the difference between the terms of the sequence
decreases as $n$ increases. Roughly speaking, then the sequence should also converge,
as the following terms don't deviate too much. This is infact the idea behind
Cauchy sequences, which allow to use restate convergence.

\begin{definition}
A Cauchy sequence is a sequence $(a_n)$ such that for any $\varepsilon > 0$, there exists an $N$ such that
$\abs{a_n - a_m} < \varepsilon$ for all $n, m \geq N$.
\end{definition}

Let us prove the proposition we just discussed, of how Cauchy sequences relate to convergence.

\begin{restatable}{theorem}{cauchyconvergence}
\label{thm:cauchy_convergence}
A sequence is cauchy if and only if it is convergent.
\end{restatable}

Unfortunately the proof of this will have to wait for a while. One direction of it is not
hard, I will leave it up to you to show that a convergent sequence is Cauchy. However,
for the other direction we'll need to cover a few more definitions. It is perhaps clear
that we would need to do this, because to prove something is convergent, you would want a value
that it converges to. For an arbitrary sequence, it seems outside our abilities
to be able to choose something like that for now.

\vspace{-1em}
\begin{exc}
\vspace{-1em}
\begin{exercise}
Show that a convergent sequence is Cauchy.
\end{exercise}
\begin{solution}
We wish to show that for any $\varepsilon > 0$, $\exists N \in \mathbb{N}$:
$\forall n,m \geq N$, $\abs{a_n - a_m} < \varepsilon$.

Since $(a_n)_{n=1}^{\infty}$ is convergent, let $N \in \mathbb{N}$ be such that
\[
\forall i > N,\quad \abs{a_i - a} < \frac{\varepsilon}{2}.
\]

Then, for any $n,m > N$,
\[
\abs{a_n - a} < \frac{\varepsilon}{2}, \qquad \abs{a_m - a} < \frac{\varepsilon}{2}
\]
\[
\implies \abs{a_n - a_m} = \abs{(a_n-a) + (a-a_m)} \leq \abs{a_n - a} + \abs{a_m - a}
< \frac{\varepsilon}{2} + \frac{\varepsilon}{2} = \varepsilon.
\]

\[
\implies \abs{a_n - a_m} < \varepsilon \text{ for any } n,m > N.
\]

Thus, $(a_n)_{n=1}^{\infty}$ is Cauchy.
\end{solution}

\begin{exercise}
Show that any Cauchy sequence is bounded
\label{exc:cauchy_bounded}
\end{exercise}

\begin{solution}
Suppose not. Note that if a sequence is not bounded, then for any $B \in
\mathbb{R}$ and any $N \in \mathbb{N}$, we can find $m \in \mathbb{N}$, $m > N$,
with $\abs{a_m} > \abs{B}$; if this were not true, the sequence would be bounded above
and below by $\max(a_1,\dots,a_N,B)$ and $\min(a_1,\dots,a_N,-B)$ respectively.

Now let $(a_n)_{n=1}^\infty$ be a Cauchy sequence, and suppose it is not
bounded. Since it is Cauchy, for any $\varepsilon > 0$, $\exists K \in \mathbb{N}$
such that
\[
\abs{a_m - a_n} < \varepsilon \quad \text{for all } m,n > K.
\]

Take $\varepsilon = 1$, and let $K \in \mathbb{N}$ be such that
$\abs{a_m - a_n} < 1$ for all $m,n > K$.

Since the sequence is not bounded, there must exist $k^* > K$ such that
\[
\abs{a_{k^*}} \geq 1 + \abs{a_K}.
\]

But then
\[
\abs{a_{k^*} - a_K} \geq \abs{a_{k^*}} - \abs{a_K} \geq 1.
\]

However, since $k^*, K > K$, $(a_n)_{n=1}^{\infty}$ being Cauchy gives
$\abs{a_{k^*} - a_K} < 1$, a contradiction.
\end{solution}
\end{exc}
\section{Limit Points}

Although a sequence might not have a limit, it might still ``cluster'' around a value
infinitely often. So we could define a new sequence, consisting of such terms in
the sequence, and it would in fact converge to such a point.
This is the idea behind \vocab{limit points}.

\begin{definition}
    We call $\theta$ a limit point of a sequence, if for all $\varepsilon > 0$, and
    $M \ge 1$, there exists $N \ge M$ such that $\abs{a_N - \theta} < \varepsilon$.
\end{definition}

Be aware of the distinction between the definition of a limit point and a limit! Here
we are just claiming we can find an arbitrarily large point close to $\theta$,
not that every point after that is also close to $\theta$. So for example, a sequence
like $a_n = (-1)^n$ has $-1$ as a limit point, but it does not converge to $-1$.

The idea of defining a new sequence derived from the original also has a name,
in particular, we define a \vocab{subsequence} of the original sequence.

\begin{definition}
	A sequence $(b_m)$ is called a \vocab{subsequence} of the original sequence $a_n$ if there
	exists a sequence $1 \le n_1 < n_2 < \dots < n_m < \dots$ such that $b_m = a_{n_m}$ for all $m$.
\end{definition}

Therefore, we should be able to find a subsequence that converges to any limit point. Since
it is a good idea to learn how to define a subsequence with specific properties,
we will prove this result.

\begin{theorem}
	If a sequence $(a_n)$ has a limit point $\theta$, then there exists a subsequence $(b_m)$
	that converges to $\theta$.
\end{theorem}

\begin{proof}
	We want to ensure that we can get arbitrarily close to $\theta$ by going
	far enough into the sequence. How do we do this? Well, the first thing is to
	utilise the definition of the limit point. We know that we can find some point, $n_1$
	that is arbitrarily close to $\theta$. But this is for a particular value of $\varepsilon$, and sure
	we could keep choosing $n_2 > n_1$ to get a term that is similarly close to $\theta$,
	but that is hardly sufficient.

	The idea is to get stricter and stricter bounds. Define a subsequence as follows, let $n_1$
	be such that $\abs{a_{n_1} - \theta} < 1$. Then let $n_2 > n_1$ be such that
	$\abs{a_{n_2} - \theta} < {1}/{2}$, and so on, that is $n_{m} > n_{m-1}$ be such that
	$\abs{a_{n_m} - \theta} < {1}/{m}$. Define $b_m = a_{n_m}$. Let us prove that this
	indeed works.

	Consider an epsilon, we must find a $N \in \NN$ such that for all $n > N$,
	$\abs{b_n - \theta} < \varepsilon$. We know that $\abs{b_n - \theta} < {1}/{n}$, so it
	suffices to choose $N$ such that ${1}/{N} < \varepsilon$, i.e. $N > {1}/{\varepsilon}$,
	which we can indeed do by the Archimedean property, so we are done.
\end{proof}

\begin{remark}
	This construction really reveals the power of sequences. You can always work with
	functions by default, but using sequences allow you to \emph{inductively} build up
	a desired sequence. This is not something you can do with functions, as might be
	obvious to you.
\end{remark}

Can a sequence cluster around any point? Clearly not if it is bounded at least. But
there is actual a stricter sense in which it can't cluster around any point---for bounded sequences---using the idea of $\limsup$ and $\liminf$.

\begin{definition}
	For a sequence $(a_n)$, the $\limsup$ is the limit of the sequence of upper bounds of
	the \emph{tails of the sequence},
	$\limsup a_n \coloneq \lim_{n \to \infty} \sup \{a_m \mid m \geq n\}$, and the $\liminf$ is the limit of the sequence of lower bounds of
	the \emph{tails of the sequence},
	$\liminf a_n \coloneq \lim_{n \to \infty} \inf \{a_m \mid m \geq n\}$.
\end{definition}

\begin{figure}[H]
\vspace{-1em}
\begin{fullwidth}
\begin{tikzpicture}
\hspace{-4em}
\begin{axis}[
    width=19cm, height=9.5cm,
    scale only axis,
    xlabel={$n$}, ylabel={$a_n$},
    xmin=0, xmax=78, ymin=-1.9, ymax=2.9,
    xtick={0,10,...,60},
    ytick={-1,-0.2,0.6,1.3,2},
    axis lines=left,
    grid=major, grid style={gray!15},
    tick label style={font=\footnotesize},
    label style={font=\small},
    title={Limit points, $\limsup$, and $\liminf$ of a bounded sequence},
    title style={font=\small, yshift=2pt},
    legend style={at={(0.5,-0.16)}, anchor=north, legend columns=3,
                  font=\footnotesize, draw=gray!40, fill=white,
                  fill opacity=0.9, text opacity=1},
    clip=false,
]

% --- shaded "forbidden past this point" bands ---
\fill[red!7]  (axis cs:0,2.1)  rectangle (axis cs:78,2.9);
\fill[blue!7] (axis cs:0,-1.9) rectangle (axis cs:78,-1.1);

% eps-boundary lines: limsup+eps and liminf-eps
\draw[dash dot, red!45!black]  (axis cs:0,2.1)  -- (axis cs:78,2.1);
\draw[dash dot, blue!45!black] (axis cs:0,-1.1) -- (axis cs:78,-1.1);

% --- the five limit points ---
\draw[densely dashed, red!55!black]  (axis cs:0,2)    -- (axis cs:62,2);
\draw[densely dotted, gray!60!black] (axis cs:0,1.3)  -- (axis cs:62,1.3);
\draw[densely dotted, gray!60!black] (axis cs:0,0.6)  -- (axis cs:62,0.6);
\draw[densely dotted, gray!60!black] (axis cs:0,-0.2) -- (axis cs:62,-0.2);
\draw[densely dashed, blue!55!black] (axis cs:0,-1)   -- (axis cs:62,-1);

\node[anchor=west, font=\footnotesize, red!55!black]  at (axis cs:65, 1.95)  {$\text{lim sup } a_n$};
\node[anchor=west, font=\footnotesize, gray!60!black] at (axis cs:65, 1.3)  {limit point};
\node[anchor=west, font=\footnotesize, gray!60!black] at (axis cs:65, 0.6)  {limit point};
\node[anchor=west, font=\footnotesize, gray!60!black] at (axis cs:65, -0.2) {limit point};
\node[anchor=west, font=\footnotesize, blue!55!black] at (axis cs:65, -0.95)  {$\text{lim inf } a_n$};

% --- callouts: nothing ever clusters beyond the eps-bands ---
\node[align=left, font=\footnotesize, text=red!55!black, draw=red!40,
      fill=red!4, rounded corners, inner sep=3pt]
  (hicall) at (axis cs:24, 2.62)
  {only $n=4$ lands above $\limsup a_i+\varepsilon$};
\draw[->, red!55!black, thin] (hicall.west) -- (axis cs:4.6,2.27);

\node[align=left, font=\footnotesize, text=blue!55!black, draw=blue!40,
      fill=blue!4, rounded corners, inner sep=3pt]
  (locall) at (axis cs:24, -1.55)
  {only $n=1,5$ land below $\liminf a_n-\varepsilon$};
\draw[->, blue!55!black, thin] (locall.west) -- (axis cs:5.4,-1.18);
\draw[->, blue!55!black, thin] (locall.west) -- (axis cs:1.4,-1.18);

% the raw sequence a_n
\addplot[only marks, mark=*, mark size=1.1pt, gray!35!black]
    coordinates { (1,-1.2000) (2,1.1000) (3,0.9667) (4,2.2500) (5,-1.2000) (6,-0.0333) (7,0.4571) (8,1.4250) (9,1.8889) (10,-0.9000) (11,-0.2909) (12,0.6833) (13,1.2231) (14,2.0714) (15,-1.0667) (16,-0.1375) (17,0.5412) (18,1.3556) (19,1.9474) (20,-0.9500) (21,-0.2476) (22,0.6455) (23,1.2565) (24,2.0417) (25,-1.0400) (26,-0.1615) (27,0.5630) (28,1.3357) (29,1.9655) (30,-0.9667) (31,-0.2323) (32,0.6312) (33,1.2697) (34,2.0294) (35,-1.0286) (36,-0.1722) (37,0.5730) (38,1.3263) (39,1.9744) (40,-0.9750) (41,-0.2244) (42,0.6238) (43,1.2767) (44,2.0227) (45,-1.0222) (46,-0.1783) (47,0.5787) (48,1.3208) (49,1.9796) (50,-0.9800) (51,-0.2196) (52,0.6192) (53,1.2811) (54,2.0185) (55,-1.0182) (56,-0.1821) (57,0.5825) (58,1.3172) (59,1.9831) (60,-0.9833) };

% sup_{n>=N} a_n  (decreasing, -> limsup)
\addplot[thick, red!55!black, mark=none]
    coordinates { (1,2.2500) (2,2.2500) (3,2.2500) (4,2.2500) (5,2.0714) (6,2.0714) (7,2.0714) (8,2.0714) (9,2.0714) (10,2.0714) (11,2.0714) (12,2.0714) (13,2.0714) (14,2.0714) (15,2.0417) (16,2.0417) (17,2.0417) (18,2.0417) (19,2.0417) (20,2.0417) (21,2.0417) (22,2.0417) (23,2.0417) (24,2.0417) (25,2.0294) (26,2.0294) (27,2.0294) (28,2.0294) (29,2.0294) (30,2.0294) (31,2.0294) (32,2.0294) (33,2.0294) (34,2.0294) (35,2.0227) (36,2.0227) (37,2.0227) (38,2.0227) (39,2.0227) (40,2.0227) (41,2.0227) (42,2.0227) (43,2.0227) (44,2.0227) (45,2.0185) (46,2.0185) (47,2.0185) (48,2.0185) (49,2.0185) (50,2.0185) };

% inf_{n>=N} a_n  (increasing, -> liminf)
\addplot[thick, blue!55!black, mark=none]
    coordinates { (1,-1.2000) (2,-1.2000) (3,-1.2000) (4,-1.2000) (5,-1.2000) (6,-1.0667) (7,-1.0667) (8,-1.0667) (9,-1.0667) (10,-1.0667) (11,-1.0667) (12,-1.0667) (13,-1.0667) (14,-1.0667) (15,-1.0667) (16,-1.0400) (17,-1.0400) (18,-1.0400) (19,-1.0400) (20,-1.0400) (21,-1.0400) (22,-1.0400) (23,-1.0400) (24,-1.0400) (25,-1.0400) (26,-1.0286) (27,-1.0286) (28,-1.0286) (29,-1.0286) (30,-1.0286) (31,-1.0286) (32,-1.0286) (33,-1.0286) (34,-1.0286) (35,-1.0286) (36,-1.0222) (37,-1.0222) (38,-1.0222) (39,-1.0222) (40,-1.0222) (41,-1.0222) (42,-1.0222) (43,-1.0222) (44,-1.0222) (45,-1.0222) (46,-1.0182) (47,-1.0182) (48,-1.0182) (49,-1.0182) (50,-1.0182) };

\end{axis}
\end{tikzpicture}
\end{fullwidth}
\vspace{-\baselineskip}\vspace{-\baselineskip}
  \sideparmargin{inner}
  \sidepar{\vspace{\baselineskip - 1em}
    \caption{\RaggedLeft The sequence $a_n = c_{n\bmod 5} + (-1)^n/n$ with $c_0,\dots,c_4 = -1,-0.2,0.6,1.3,2$.
    One can see that
    five subsequences cluster at five distinct values.
    Red/blue curves are $\sup_{n\ge N}a_n$ and $\inf_{n\ge N}a_n$, which converge to $\limsup a_i$
    and $\liminf a_n$. For any $\varepsilon>0$ (here $\varepsilon=0.1$),
    only finitely many terms lie above $\limsup a_i+\varepsilon$ or below
    $\liminf a_n-\varepsilon$.}}
\end{figure}

We will show that these exist for convergent sequences (more generally, for bounded sequences), and these are the tools
that will finally allow us to show that all Cauchy sequences are convergent. But
first let us discuss about the idea of clustering I mentioned---what the $\limsup$ and $\liminf$ are really doing.
As I mentioned, they are the greatest and smallest limit points. That is, the function clusters around
a value greater than the $\limsup$ and smaller than the $\liminf$ only finitely many times.

Before we prove \Cref{thm:cauchy_convergence}, let us first discuss this a little. We
want to show that every bounded sequence has $\limsup$ and $\liminf$ as limit points. But
to do that at all, we must show these actually exist, based on our definition.

\begin{theorem}
	If $(a_n)$ is a bounded sequence, then the sequence defined by $t_n = \sup\{a_m \mid m \ge n\}$
	and $s_n = \inf\{a_m \mid m \ge n\}$ is convergent, and $\limsup a_n = \lim_{n\to\infty} t_n$ and $\liminf a_n = \lim_{n\to\infty} s_n$.
	\label{prop:limsup_liminf_exist}
\end{theorem}

\begin{proof}
I will only do one of these, and leave the other for you to fill in.

We will do this in a silly manner, first showing that $(t_n)$ is bounded,
and then that is monotonically decreasing\footnote[0cm]{The definition is analogous to that of
monotonically increasing sequence, i.e $t_n \ge t_m$ for all $m \ge n$.}. This will allow
us to use \Cref{exc:monotone_convergence} to conclude that $(t_n)$ is convergent.

The first is easy to conclude from \Cref{exc:cauchy_bounded}. Since $(a_n)$ is bounded,
there exists $M$ such that $\abs{a_n} \le M$ for all $n$. Then,
$t_n \ge a_n \ge -M$ for all $n$. Also, since $M$ is an upper bound for $\{a_m \mid m \ge n\}$,
$t_n \le M$ for all $n$. Therefore, $(t_n)$ is also bounded.

\marginnote{Note that the $\sup$ exists because the sequence is bounded, so the
tail has an upper bound.}

Now for the latter, we note that $t_n$ is the least upper bound of $\{a_m \mid m \ge n\}$,
and $t_{n+1}$ is the least upper bound of $\{a_m \mid m \ge n+1\}$. If $t_n < t_{n+1}$,
then $t_{n+1}$ is not the least upper bound of $\{a_m \mid m \ge n + 1\}$, $t_n \ge a_m$ for all $m \ge n + 1$,
so it is an upper bound, and $t_n < t_{n+1}$, which is a contradiction. Therefore, we must have that
$t_n \ge t_{n+1}$, and so $(t_n)$ is monotonically decreasing.

To prove that it is convergent, we will use \Cref{exc:monotone_convergence} with a bit of trickery.
It is easy to see that $t'_n = -t_n$ is bounded and monotonically increasing, so it
must be convergent. Also, $-\mathbb{1}_n = -1$ is also convergent (it converges to $-1$), so
$t_n = -\mathbb{1}_n \cdot t'_n$ must also be convergent by \Cref{exc:limit_properties}.
By definition, we write $\limsup a_n = \lim_{n \to \infty} t_n$.
\end{proof}

Now let us consider the result we talked about:

\begin{theorem}
    If $(a_n)$ is a bounded sequence, then $\limsup a_n$ is a limit point of $(a_n)$.
\end{theorem}

\begin{proof}
We must show that for all $\varepsilon > 0$ and $M \ge 1$, there exists $N \ge M$ such that
    $\abs{a_N - \limsup a_i} < \varepsilon$.

    Fix $\varepsilon > 0$ and $M \ge 1$. Since $t_n \to \limsup a_i$, choose $N_1$ such that
    $\abs{t_n - \limsup a_i} < \varepsilon/2$ for all $n \ge N_1$. Let $n_0 = \max(N_1, M)$, so that
    $\abs{t_{n_0} - \limsup a_i} < \varepsilon/2$ and $n_0 \ge M$.

    By definition, $t_{n_0} = \sup\{a_m \mid m \ge n_0\}$ is the \emph{least} upper
    bound of that set. If every $a_m$ with $m \ge n_0$ satisfied
    $a_m \le t_{n_0} - \varepsilon/2$, then $t_{n_0} - \varepsilon/2$ would itself be an
    upper bound smaller than $t_{n_0}$, contradicting leastness. So there exists some
    $N \ge n_0 \ge M$ with
    \[
        a_N > t_{n_0} - \varepsilon/2.
    \]
    Also $a_N \le t_{n_0}$, since $t_{n_0}$ is an upper bound for $\{a_m \mid m \ge n_0\}$
    and $N \ge n_0$. Combining,
    \[
        \abs{a_N - t_{n_0}} < \varepsilon/2.
    \]
    Now use the triangle inequality together with $\abs{t_{n_0} - \limsup a_i} < \varepsilon/2$:
    \[
        \abs{a_N - \limsup a_i} \le \abs{a_N - t_{n_0}} + \abs{t_{n_0} - \limsup a_i} < \varepsilon/2 + \varepsilon/2 = \varepsilon.
    \]
    Since $N \ge M$, this is exactly what was required. As $\varepsilon > 0$ and $M \ge 1$
    were arbitrary, $\limsup a_i$ is a limit point of $(a_n)$.
\end{proof}

\marginnote{I write $\limsup a_i$ instead of $\limsup a_n$ to avoid notational overload. The index
doesn't really matter.}[-8cm]

Let us now prove \Cref{thm:cauchy_convergence}.
\cauchyconvergence*

\begin{proof}
Let $a_n$ be a Cauchy sequence. We claim that it converges to $\limsup a_i$. This exists because
of \Cref{exc:cauchy_bounded} and \Cref{prop:limsup_liminf_exist}.

We wish to find a $N$ such that $\abs{a_n - \limsup a_i} < \varepsilon$ for all $n \ge N$.
Since $(a_n)$ is a Cauchy sequence, there exists a $N_1$ such that $\abs{a_n - a_m} < \varepsilon/4$ for all $n, m \ge N_1$.

Since $t_n \coloneq \sup \{a_m \mid m \ge n\}$ is convergent, there exist $N_2$
such that $\abs{t_n - \limsup a_i} < \varepsilon/2$ for all $n \ge N_2$, and so for $n > N_3 = \max(N_1, N_2)$.

Now consider the terms $a_{N_3}, a_{N_3+1}, \dots$. Note that $t_n $ is greater than all of them
by definition. Now there must exist a $i$ such that $t_n - a_{N_3 + i} < \varepsilon/4$, otherwise
$t_n - \varepsilon/4$ would be a smaller upper bound of $\{a_m \mid m \ge N_3 + i\}$, which is a contradiction.

Let this be $a_{N}$. Since $N > N_3 = \max(N_1, N_2)$, we have that $\abs{a_N - a_m} < \varepsilon/4$ for all $m \ge N$.
Using this, and $t_n - a_{N} = \abs{t_n - a_{N}} < \varepsilon/4$, we get that:
\[
\varepsilon/2 > \abs{a_{N} - a_m} + \abs{t_n - a_{N}} \ge \abs{a_m - t_n}
\]
and finally, using $\abs{t_n - \limsup a_i} < \varepsilon/2$, we get that:
\[
\varepsilon > \abs{t_n - \limsup a_i} + \abs{a_m - t_n} \ge \abs{a_m - \limsup a_i}
\]
for all $m \ge N$. Thus we are done.
\end{proof}

\vspace{-1em}
\begin{remark}
    Continuing on \Cref{rem:extended_convergence}, we may also
    use the extended definition to allow $\limsup a_n$ to be $\infty$ and
    $\liminf a_n$ to be $-\infty$. This is a surprisingly important idea,
    as it allows us to say that \emph{all} sequences, regardless of them being bounded or not,
    have a $\limsup$ and $\liminf$.
\end{remark}
\vspace{-1em}
\begin{remark}
    Here is a trick I have pulled over you to prove \Cref{thm:cauchy_convergence}. Cauchy
    sequences are what we use to \emph{construct} the reals. Look carefully at what properties
    I used for the proof---the most important is the existence of $\sup$. In particular,
    Cauchy sequences are not all convergent if we work with the rationals! We construct
    the reals by defining a real to be an equivalence class of Cauchy sequences whose
    difference is also a Cauchy sequence (i.e the difference between their terms decreases to zero as
    we grow larger).
\end{remark}

\begin{exc}
    \begin{exercise}[subtitle = {Balzano Weirstrass Theorem}]

    Prove that every bounded sequence has a convergent subsequence.

    \emph{Hint: You already know one such subsequence! What do subsequences of bounded sequences converge to?
        But I also want you to do this in another manner, without appealing to what we have
        already shown, i.e without using the theorems we have already proved for bounded
        sequences.}
    \end{exercise}
    \begin{solution}
    \begin{claim}
    Any sequence $(a_n)$ contains a subsequence which is either nondecreasing (montonically
    increasing) or non-increasing (monotonically decreasing).
    \end{claim}
    \begin{proof}
    Call a natural number $n$ a \emph{peak} point of the sequence $(a_n)$ if $a_m < a_n$ for all
    $m > n$

    Case 1. The sequence has infinitely many peak points. In this case, if $n_1 < n_2 <
    n_3 < \ldots$ are the peak points, then $a_n_1 > a_n_2 > a_n_3 > \ldots$ , so $a_n_k$ is the desired
    (nonincreasing) subsequence.

    Case 2. The sequence has only finitely many peak points. In this case, let $n_1$ be greater
    than all peak points. Since $n_1$ is not a peak point, there is some $n_2 > n_1$ such that
    $a_n_2 \geq a_n_1$ . Since $n_2$ is not a peak point (it is greater than $n_1$ , and hence greater
    than all peak points) there is some $n_3 > n_2$ such that $a_n_3 \geq a_n_2$ . Continuing in this
    way we obtain the desired (nondecreasing) subsequence.

    Since we have a monotonically increasing/decreasing sequence, and it must bounded
    because it is a subsequence of a bounded sequence, by \Cref{exc:monotone_convergence}
    we are done.
    \end{proof}
    \end{solution}
    \begin{exercise}
    Let $(a_n)$ be a bounded sequence. Show that:
    \begin{enumerate}
    \item $\liminf a_n$ exists.
    \item $\liminf a_n \leq \limsup a_n$.
    \item $\liminf a_n = \limsup a_n$ if and only if $(a_n)$ is convergent (or Cauchy).
    \end{enumerate}
    \end{exercise}
\end{exc}
